\section{Hardware-Sovereign Serving}
\label{sec:sovereign}

Serving an open-weight LLM is, in practice, often equivalent to depending on a
single vendor's stack (CUDA/NVIDIA). For a system whose remit is national law,
that coupling is itself a sovereignty risk. Our serving layer is a thin
\emph{OpenAI-compatible} indirection: switching silicon is switching three
environment variables (\texttt{LLM\_BASE\_URL}, \texttt{LLM\_API\_KEY},
\texttt{LLM\_MODEL}); the pipeline is unchanged.

Table~\ref{tab:backends} states, honestly, what runs where. We demonstrate the
pipeline \emph{live} on a Qualcomm Cloud~AI~100 inference
accelerator~\cite{qualcommcloudai100} (Llama-3.1-8B), a dedicated
inference silicon whose value proposition is operations-per-watt rather than
raw training throughput. We target AMD via a portable inference
runtime~\cite{modularmax2024}; European accelerators are stated as vision, not
measurement. We report real latency and tokens/s; we deliberately do
\emph{not} report a fabricated energy figure, citing only that dedicated
inference silicon is designed for high throughput-per-watt.

\begin{table}[t]
\centering
\caption{Serving backends and honest status. The pipeline is identical across
rows; only the OpenAI-compatible endpoint changes.}
\label{tab:backends}
\small
\begin{tabular}{@{}lll@{}}
\toprule
Backend & Silicon & Status \\
\midrule
Qualcomm Cloud AI 100 & Qualcomm & \textbf{live, proven} \\
Modular MAX & AMD Instinct & target (runtime supports it) \\
Modular MAX & NVIDIA & backward compatibility \\
Self-hosted (open weights) & NVIDIA (cloud) & fallback \\
VSora / Axelera & EU/FR & vision (no access) \\
\bottomrule
\end{tabular}
\end{table}
